% !TeX root = ../papers/nc-p3.tex
\section{共通二次超曲面の既約軌跡と二つの中間座標}
\label{sec:trisecant-common-quadric-locus}
\begin{proposition}\label{prop:trisecant-common-quadric-locus}
$R\subset\mathbb P^4$ を有理正規四次曲線に固定する。
$R$ と $T^{-\mathsf t}R$ が共通二次式を持つ $T\in\mathrm{PGL}_5$ の軌跡は
余次元4の既約閉集合である。一般の共通二次超曲面は滑らかで一意である。
\end{proposition}
\begin{proof}
$R$ の二次イデアルは六次元である。これを移動した曲線の八次切断九空間へ制限する
行列の最大六次小行列式が条件を与え、空でない各成分の余次元は高々4である。
以下、共通二次式を印した入射を取る。$R$ の二次式の射影五空間の中の二点 $Q,Q'$ を選び、
それらを運ぶ射影変換を選ぶ表示になる。どの二次式も階数は3以上である。
階数5の底は既約で次元10、上のファイバーは次元10の連結群
$\mathrm{PGO}_5$ のトーサーである。全空間は既約となる。
実際滑らかな二次形式の基底をエタール局所的に選べるので射は滑らかで、
幾何学的一般ファイバーが既約であり、各成分が底を支配する。

$R$ 上消える階数4以下の二次式は次元4以下である。行列式は恒等零でなく、
例えば逆対角成分 $(1,3,-8,3,1)$ の行列が可逆である。
階数3のものは次元2以下である。階数3の式を $LM-H^2$ と書き、
$L,M,H$ を $R$ 上の二項四次式と見る。これは一次形式全体と $S_4$ との同型である。
一意分解から $L=d a^2,H=d ab,M=d b^2$ とでき、互いに素な $a,b$ の
共通次数は1か2である。式は $[d]$ と $\langle a,b\rangle$ で決まり、
次数1なら二次式 $[d]$ の二次元、次数2なら $\operatorname{Gr}(2,S_2)$ の二次元である。
対称五行行列の階数 $r$ 軌道の余次元は $(5-r)(6-r)/2$ である
（根基を選んで商上の非退化対称形式を選ぶ）。
従って $r=4,3$ の安定群次元は $11,13$、対応する入射の次元は
それぞれ $2\cdot4+11=19$、$2\cdot2+13=17$ 以下である。
軌跡の各成分は次元20以上なので、全て階数5入射の像の閉包に入り、
その次元は20である。一般ファイバーは零次元で、共通二次式は射影線型空間を
なすからただ一点である。
\end{proof}

\begin{lemma}\label{lem:two-middle-rigidity}
四辺有限の標準有限Frobenius候補の標識付き配置を $(X,C,Y,G)$ とし、
$C,G$ が標識付き射影非同値とする。$C$ から $G$ への対称二次式行列で
階数2以下の軌跡が有限なら、同じ標識付き反対対 $(X,Y)$ を持つ四辺有限の標準正則窓は
この候補か、その順序反転と基底変更で一致する。
\end{lemma}
\begin{proof}
他の奇位置 $V$ も $X$ との有限完全交叉の跡により $C^2$ と単元倍まで同じ平方を持つ。
定理\ref{thm:double-point-opposite-reconstruction}と非被約での階数2排除
（定理\ref{thm:nonreduced-opposite-root-separation}）によりそのVeronese像は
$C$ 自身または $G$ のものに等しい。Veronese同型が各印付き評価ベクトルを対応させるため
座標配置も射影同値である。従って三つの連続座標は $(X,C,Y)$ または $(X,G,Y)$ となる。
各隣接積の核が関係六空間で、三次交叉は有限Frobenius条件でちょうど四空間である。
その四成分を最後の文字に束ねれば四テンソルはその文字の基底を除き一意である。
後者の三座標は最初の候補を反転すれば得られる。添字保持同型で反転を同一視する必要はない。
\end{proof}

\begin{theorem}\label{thm:general-common-quadric-trisecant-exclusion}
命題\ref{prop:trisecant-common-quadric-locus}の軌跡の稠密開集合では、
標準座標以外を許しても、その標識付き三割平面の反対対を含む四辺有限な標準正則延長はない。
\end{theorem}
\begin{proof}
共通二次式が一意な開集合で $\mathcal G$ は階数4の束で、前節の四テンソルが族として定まる。
その二つの隣接切断の階数は6以下である。$x,C$ と $C,y$ の六関係空間に
テンソルが入ること、および対向する辺が $S_5$ で収縮表示されることによる。
前節の有理標本で両階数6、反対階数16、有限Frobeniusと全射影の有限性が成り立つ。
従って開近傍でも成立する。特に三次、四次の交叉次元は収縮により下から抑えられ、
上からの開条件で標本と一致する。
$C,G$ が標識付き射影同値なら、$C,y,G$ の二辺が反転となる。
命題\ref{prop:finite-flip-obstructions}の三次交叉の対称性により四テンソルは
奇位置で対称となり、反対切断階数16に反する。この対称性は有限完全交叉のKoszul分解だけを使う。
従って有限Frobenius窓の定理\ref{thm:all-finite-opposite-octics}により
$C$ から $G$ への二次式行列の行列式は非零で、$C$ の二十点で三重に消える。

\path{computations/verify_trisecant_common_quadric_locus.py} は同じ標本で
$C_{00}$ 倍の可逆性、$C_{01}/C_{00}$ の特性多項式の平方自由性を良い素数1009で確認する。
従って標数零でも二十点は被約である。二階偏導を四次最大小行列式の六次商へ送る
$200\times165$ 行列の階数は164であり、核の八次式を
$C_{11}=C_{00}+2C_{01}+5C_{10}$ に制限すると三偏導の十二次倍が十九次式全体を張る。
標数零で既に行列式が核にあるから一次元核はこの良い還元に持ち上がる。
一般の $T$ でも被約性、核の一意性、平面切断の滑らかさは保たれる。
従って八次曲面の特異点は有限で、特に二次式行列の階数2以下は有限である。
補題\ref{lem:two-middle-rigidity}で可能な二つの順序はいずれも
定理\ref{thm:trisecant-middle-simple-obstruction}で排除される。
反転には左右の正則性の同値を用いる。
\end{proof}
